% ============================================================
% Parecer da banca Banca_TCC sobre Artigo.tex (RadarTorres)
% Gerado pelo agente banca-tcc (.claude/agents/banca-tcc.md).
% Não editar à mão sem regenerar a partir de uma nova análise.
% Compile com pdfLaTeX.
%
% QUARTA revisão. As anteriores (27/10, nota 6,33; segunda revisão, nota
% 6,20; terceira revisão, nota 7,9) estão em Parecer_Banca_TCC.tex e
% Parecer_Banca_TCC_rev3.tex e permanecem intactas. Esta verifica duas
% mudanças relatadas desde a 3ª revisão -- remoção dos genes treinados do
% versionamento e reescrita da frase da Seção 2.3.5 -- e encontra que
% nenhuma das duas alcança a branch que parece ser a de entrega.
% ============================================================
\documentclass[11pt,a4paper]{article}

\usepackage[T1]{fontenc}
\usepackage[utf8]{inputenc}
\usepackage[brazil]{babel}
\usepackage[a4paper,margin=2.5cm]{geometry}
\usepackage{xcolor}
\usepackage{longtable}
\usepackage{array}
\usepackage{url}

\title{\textbf{Parecer da Banca Examinadora --- TCC ``RadarTorres''}\\[0.3em]
\large Quarta revisão --- gerado pelo agente \texttt{banca-tcc}}
\author{}
\date{Análise referente a \texttt{Artigo.tex}/\texttt{Artigo.pdf} (RadarTorres, TCC de Engenharia da
Computação, UCL)\\[0.3em]
Auditado contra o código-fonte, os artefatos em JSON, o banco \texttt{radartorres.db} em uso e o
histórico git de \textbf{duas} branches: \texttt{main} e \texttt{banco-dados}}

\newcommand{\critico}{\textbf{\color{red!60!black}[CRÍTICO]}}
\newcommand{\alto}{\textbf{\color{orange!70!black}[ALTO]}}
\newcommand{\moderado}{\textbf{\color{blue!50!black}[MODERADO]}}
\newcommand{\baixo}{\textbf{\color{gray}[BAIXO]}}
\newcommand{\retirado}{\textbf{[RETIRADO]}}

\begin{document}
\maketitle

% ============================================================
% QUADRO-RESUMO
% ============================================================
\begin{center}
\fcolorbox{black}{gray!8}{%
\begin{minipage}{0.93\textwidth}
\vspace{0.4em}
\begin{center}
{\Large\textbf{Nota: 7,7 / 10}}\\[0.2em]
{\small pareceres anteriores: 6,33 (1ª revisão) \quad$\cdot$\quad 6,20 (2ª revisão) \quad$\cdot$\quad
7,9 (3ª revisão)}
\end{center}
\vspace{0.3em}
\hrule
\vspace{0.6em}
\textbf{Veredicto.} As duas correções que a 3ª revisão pedia como bloqueantes foram feitas ---
\textbf{mas na branch \texttt{main}}, enquanto a branch que tudo indica ser a de entrega
(\texttt{banco-dados}, cujo último commit se chama literalmente \textit{``revisão do Gabriel no dia
da entrega''}, e é a que está hoje fisicamente aberta neste repositório) \textbf{não recebeu nenhuma
das duas}. Nela, \texttt{rede-controlador.net} e \texttt{genoma-controlador.json} continuam
rastreados pelo git e embutidos pelo \texttt{.csproj}, e a Seção 2.3.5 continua com a frase antiga
--- ``o repositório não versiona uma rede treinada'' ---, que volta a ser falsa pelos mesmos motivos
da 1ª e da 2ª revisão.

\vspace{0.5em}
\textbf{Pendência que muda a nota: 1, mas é nova e mais grave que a antiga.}
Não é mais ``reescrever uma frase'': é reconciliar duas branches divergentes antes de decidir qual
delas vai a exame --- a que tem a correção não tem o conteúdo que esta própria banca elogiou (a
explicação da tanh, a dívida técnica dos ciclos de dependência, o item (vii) de trabalhos futuros), e
a que tem o conteúdo não tem a correção.

\vspace{0.5em}
\textbf{Achados:} 1 crítico reaberto \quad$\cdot$\quad 2 altos \quad$\cdot$\quad 3 moderados
\quad$\cdot$\quad 3 baixos \quad$\cdot$\quad 1 retirado
\vspace{0.6em}
\end{minipage}}
\end{center}

\vspace{0.8em}

\textit{Como ler este parecer:} a Seção A resume; a B justifica a nota critério por critério; a C
compara com o parecer anterior; a D lista os problemas, cada um com gravidade no título; a E mostra o
que foi conferido contra o código e os dados, com veredicto por item; a I traz as pendências em
ordem de prioridade; a J diz o que \emph{não} foi verificado. Quem tem pouco tempo: leia A, H e I.

\vspace{0.5em}

\section*{A. Síntese}

Esta rodada foi aberta para verificar duas afirmações: que os genes treinados saíram do
versionamento da \texttt{main} e foram preservados em \texttt{genes-treinamento}, e que a frase da
Seção 2.3.5 foi reescrita para não afirmar isso que não é mais verdade. As duas são verdadeiras ---
\textbf{na branch \texttt{main}}. \texttt{git log --all} mostra os três commits citados
(\texttt{2863fc4}, \texttt{0f58ccb}, \texttt{501ce86}), todos alcançáveis a partir de \texttt{main},
e \texttt{git show 501ce86} confirma que a nova frase é exatamente a que a 3ª revisão recomendou.

O problema é que \texttt{main} não é a branch que este repositório está mostrando como estado atual
do trabalho. O diretório de trabalho está em \texttt{banco-dados}, cujo commit de ponta --- feito às
14:26 do mesmo dia, \textbf{antes} dos dois commits de correção em \texttt{main} (15:03 e 15:40) --- se
chama \texttt{a8a7b16 ``revisão do Gabriel no dia da entrega''}. \texttt{main} e \texttt{banco-dados}
divergem desde o commit \texttt{b0af938}: nenhuma é ancestral da outra, e nenhum merge posterior as
reconcilia. Nessa branch, \texttt{git ls-tree} mostra os dois arquivos --- \texttt{rede-controlador.net}
e \texttt{genoma-controlador.json} --- ainda rastreados; o \texttt{.gitignore} dessa branch não tem a
linha que os exclui; e a Seção 2.3.5 do \texttt{Artigo.tex} efetivamente presente no disco, nesta
branch, ainda diz ``o repositório não versiona uma rede treinada: um executável compilado sem
treinamento prévio na máquina usa a heurística pura'' --- a frase que a 1ª, a 2ª e a 3ª revisão
apontaram como falsa, porque o \texttt{.csproj} embute o arquivo sob
\texttt{Condition="Exists(...)"}, a condição está satisfeita, e o arquivo está no disco.

Em outras palavras: \textbf{o achado crítico das três revisões anteriores não foi corrigido onde
importa} --- ele foi corrigido em uma branch e documentado como corrigido tomando por base outra. E a
branch que recebeu a correção, \texttt{main}, por sua vez \textbf{não contém} parte do que esta banca
já havia registrado como melhoria (a explicação didática da tanh, a dívida técnica dos ciclos de
dependência, o item (vii) de trabalhos futuros, o rótulo \texttt{sec:arquitetura}) --- conteúdo que
só existe em \texttt{banco-dados}. As duas branches têm, cada uma, metade do que o artigo precisa ser.

Os demais oito achados da 3ª revisão (D.2, D.4 e D.6 a D.10) foram reconferidos nesta rodada contra o
\texttt{Artigo.tex} tal como está no disco hoje (branch \texttt{banco-dados}) e persistem
\textbf{idênticos}, byte a byte nas frases citadas. Nenhum foi corrigido; nenhum piorou. A Tabela 1,
a tabela central de resultados, não foi reafetada por nada desta rodada e continua reproduzindo
exatamente.

\textbf{Saldo.} O método e os resultados centrais continuam sólidos e verificáveis --- isso não mudou.
O que mudou é que o processo de correção criou um problema novo, de natureza diferente do antigo:
antes havia uma frase desalinhada do repositório; agora há dois repositórios parcialmente corrigidos
e nenhum critério registrado de qual deles é o artigo. É uma falha de processo de entrega, não de
método de pesquisa, e tem conserto barato --- mas, enquanto não for resolvida, reabre exatamente o
tipo de achado que bloqueia defesa: uma afirmação central do texto (``o repositório não versiona uma
rede treinada'') que é falsa na branch que está fisicamente aberta e pronta para ser mostrada à banca.

\section*{B. Nota geral}

\textbf{Nota ponderada: $\approx 7{,}70\,/\,10$} \quad (pareceres anteriores: $6{,}33$, $6{,}20$ e
$7{,}9$).

\textit{A nota caiu 0,2 em relação à 3ª revisão, não porque algo piorou no método, mas porque o
crédito que a 3ª revisão deu a Reprodutibilidade (de 6,0 para 7,0) foi condicionado à correção estar
``conferida em \texttt{origin/main}'' --- o que era verdade e continua sendo, mas revelou-se
insuficiente, porque \texttt{origin/main} não é o estado do repositório que a branch de trabalho
mostra. Reprodutibilidade volta a 5,5: mais baixa que antes da correção (6,0), porque o problema de
hoje é mais grave que o original --- não é mais ``uma frase errada'', é ``duas branches, cada uma
certa a seu próprio modo e nenhuma delas completa''. Os outros cinco critérios não mudam: nenhum dos
oito achados persistentes (D.2, D.4, D.6--D.10) foi afetado por esta rodada.}

\renewcommand{\arraystretch}{1.3}
\begin{longtable}{|p{4.1cm}|p{1.1cm}|p{1.1cm}|p{1.4cm}|p{6.6cm}|}
\hline
\textbf{Critério} & \textbf{Peso} & \textbf{Nota} & \textbf{Contrib.} & \textbf{Justificativa (resumo)} \\
\hline
\endhead
Fundamentação teórica & 20\% & 8,0 & 1,600 &
Inalterada. As 21 referências e a disciplina de amarrar cada uma a uma decisão de projeto persistem
na branch auditada (\texttt{banco-dados}); nada nesta rodada afeta este critério. \\
\hline
Metodologia & 20\% & 8,0 & 1,600 &
Inalterada. Seeds reservadas, protocolo de cinco execuções, Wilcoxon pareado: tudo conferido na 3ª
revisão e não reafetado. Ainda puxada por D.2 (população ``padrão'' já não é 60) e D.7 (teto de
alvos por quadrante inexistente), ambos persistentes e reconferidos nesta rodada. \\
\hline
Resultados e análise & 25\% & 7,5 & 1,875 &
Inalterada. A Tabela 1 continua reproduzindo exatamente; D.5 (Tabela 2 sem artefato) e D.6 (recompensa
por abate não é fixa em 10) persistem sem mudança. \\
\hline
Redação e estrutura & 15\% & 8,5 & 1,275 &
Inalterada quanto ao conteúdo da branch auditada. O \texttt{.log} de 27/11 14:25 (branch
\texttt{banco-dados}) não cobre o estado de \texttt{main} e por isso não foi usado para julgar a
versão corrigida; a aritmética de D.4 persiste. \\
\hline
Contribuição e originalidade & 10\% & 8,0 & 0,800 &
Inalterada. O resultado negativo continua declarado no Resumo, no Abstract e na Conclusão da branch
auditada; nenhuma mudança de conteúdo nesta rodada. \\
\hline
Reprodutibilidade & 10\% & 5,5 & 0,550 &
\textbf{Desce de 7,0.} A correção de D.1 existe, mas em \texttt{main}, branch que não está
fisicamente aberta nem é ancestral/descendente de \texttt{banco-dados}. Na branch auditada, o genoma
continua versionado e embutido, e a frase continua falsa --- o achado crítico original está
\textbf{integralmente reaberto}. Mais grave que a nota de 6,0 anterior à 3ª revisão, porque agora há
também confusão de qual branch representa o trabalho a ser defendido. \\
\hline
\multicolumn{3}{|r|}{\textbf{TOTAL}} & \textbf{7,70} & \\
\hline
\end{longtable}

\subsection*{Por que a nota caiu, apesar de o método não ter mudado}

Porque a gravidade de um achado é definida pelo impacto na conclusão lida pela banca no exemplar que
tem na mão, não pela existência de uma correção em algum lugar do histórico git. Se a banca abrir
\texttt{banco-dados} --- o que ela fará, porque é a branch que está aberta --- vai ler exatamente a
mesma frase falsa que a 1ª revisão apontou há três rodadas. A correção feita em \texttt{main} é real,
está bem feita e usa a redação certa; ela só não chega a quem vai ler o artigo entregue, a menos que
alguém funda as duas branches antes da defesa.

\section*{C. Evolução desde o parecer anterior}

\subsection*{C.1 Corrigido ou materialmente melhorado}

Nenhum item. As duas correções relatadas para esta rodada (remoção dos genes do versionamento;
reescrita da frase da Seção 2.3.5) foram de fato feitas, com qualidade, e teriam resolvido D.1 por
completo --- \textbf{mas na branch \texttt{main}}, que não é a branch auditada na 3ª revisão nem a
que está hoje aberta neste diretório. Até que as duas branches sejam reconciliadas, nenhum problema
pendente da 3ª revisão pode ser dado como corrigido no artigo que será efetivamente apresentado.

\subsection*{C.2 Melhorias novas, não pedidas no parecer anterior}

Nenhuma. A branch \texttt{banco-dados} não recebeu nenhum commit de conteúdo do artigo desde a 3ª
revisão; o único commit novo nela é o próprio parecer anterior, que ainda está sem commit
(\texttt{Parecer\_Banca\_TCC\_rev3.tex} e \texttt{.pdf} aparecem como \textit{untracked} no
\texttt{git status} desta auditoria).

\subsection*{C.3 Persiste, idêntico}

\begin{enumerate}
  \item \textbf{D.2} (população ``padrão do sistema'' já não é 60) --- frase idêntica, linha 235 do
  \texttt{Artigo.tex} atual.
  \item \textbf{D.4} (partição aritmética ``sete cadastrais e quatro de treinamento'' vs.\ tabela
  8+3) --- idêntica, linha 522.
  \item \textbf{D.5} (Tabela 2 sem artefato congelado: 134,17 / 119,62 / 137,35 não existem no banco
  em uso) --- idêntica, linhas 643--644.
  \item \textbf{D.6} (recompensa por abate descrita como 10 pontos fixos; o código soma 7 a 10) ---
  idêntica, linha 470.
  \item \textbf{D.7} (``cada quadrante admite até cinco alvos simultâneos''; não existe teto no
  código) --- idêntica, linha 469.
  \item \textbf{D.8} (``4631 oportunidades'' sem indicar que é por execução) --- idêntica, linha 607.
  \item \textbf{D.9} (execuções reportadas a partir de uma cópia em \texttt{C:\textbackslash TEMP}) ---
  não reconferida nesta rodada por orçamento; sem motivo para supor mudança.
  \item \textbf{D.10} (especificações de hardware autodeclaradas) --- sem mudança.
  \item \textbf{O problema de fundo de D.1}: genoma versionado, embutido por \texttt{Exists(...)} no
  \texttt{.csproj}, divergente dos dez campeões relatados. \textbf{Persiste, idêntico, na branch
  \texttt{banco-dados}} --- apesar de corrigido em \texttt{main}.
\end{enumerate}

\subsection*{C.4 Piorou ou surgiu}

\begin{enumerate}
  \item \textbf{Bifurcação não reconciliada entre \texttt{main} e \texttt{banco-dados}.} As duas
  branches divergem desde \texttt{b0af938}; nenhuma é ancestral da outra
  (\texttt{git merge-base --is-ancestor} falha nos dois sentidos). \texttt{main} ganhou os dois
  commits de correção de D.1 (\texttt{0f58ccb}, \texttt{501ce86}); \texttt{banco-dados} ganhou dois
  commits próprios (\texttt{cc11e44}, \texttt{a8a7b16 ``revisão do Gabriel no dia da entrega''}) que
  não tocam D.1. Nenhum commit de merge posterior aparece em nenhuma das duas branches locais ou
  remotas. \textbf{Isto é novo nesta rodada} --- não existia como achado nas revisões anteriores,
  porque até a 3ª revisão as duas branches ainda não tinham se afastado neste ponto específico do
  artigo.
  \item \textbf{Consequência prática:} a correção do achado mais grave das três revisões anteriores
  (D.1) não produz efeito na branch que parece ser a de entrega. O risco concreto é o autor acreditar
  --- com razão, se olhar só \texttt{main} --- que o item está resolvido, e a banca real abrir
  \texttt{banco-dados} e encontrar a frase falsa e o arquivo versionado, exatamente como na 1ª
  revisão.
\end{enumerate}

\section*{D. Problemas}

\renewcommand{\arraystretch}{1.25}
\begin{longtable}{|p{0.9cm}|p{2.6cm}|p{11.6cm}|}
\hline
\textbf{Item} & \textbf{Gravidade} & \textbf{Em uma linha} \\
\hline
\endhead
D.1 & \critico\ (reaberto) & A correção existe só em \texttt{main}; na branch aberta
(\texttt{banco-dados}) o genoma continua versionado/embutido e a frase continua falsa \\
\hline
D.2 & \alto & ``População de 60, o padrão do sistema'': o padrão hoje é 100, mudou depois do
experimento --- persiste \\
\hline
D.3 & \retirado & Contagem de tabelas do banco --- argumento de escopo dos autores aceito na 3ª
revisão; mantido retirado \\
\hline
D.4 & \baixo & ``Sete cadastrais e quatro de treinamento'' não bate com as 8+3 linhas da própria
tabela --- persiste \\
\hline
D.5 & \alto & A Tabela 2 não é reproduzível: o banco que a originou foi sobrescrito --- persiste \\
\hline
D.6 & \moderado & ``10 pontos por alvo interceptado'': o código soma de 7 a 10, conforme a distância
--- persiste \\
\hline
D.7 & \moderado & ``Cada quadrante admite até cinco alvos simultâneos'': não existe esse teto no
código --- persiste \\
\hline
D.8 & \baixo & ``4631 oportunidades de acionamento'' é por execução, não o total --- persiste \\
\hline
D.9 & \baixo & As execuções reportadas partiram de uma cópia em \texttt{C:\textbackslash TEMP} ---
não reconferido, sem indício de mudança \\
\hline
D.10 & \baixo & Especificações de hardware são autodeclaradas --- persiste \\
\hline
\end{longtable}

\subsection*{D.1 \critico\ (reaberto) A correção de D.1 existe --- na branch errada}

\textbf{O que foi anunciado para esta rodada.} Que \texttt{rede-controlador.net} e
\texttt{genoma-controlador.json} foram removidos do versionamento da \texttt{main} (\texttt{git rm
--cached}, commit \texttt{0f58ccb}) e preservados na branch \texttt{genes-treinamento}
(\texttt{2863fc4}), com \texttt{.gitignore} atualizado; e que a frase da Seção 2.3.5 foi reescrita
para ``a versão entregue não inclui rede treinada; um executável compilado a partir dela usa a
heurística pura, comportamento de fábrica previsto no projeto'' (commit \texttt{501ce86}).

\textbf{O que esta auditoria confirma sobre \texttt{main}.} As duas afirmações são verdadeiras.
\texttt{git log --all} mostra os três commits, todos alcançáveis a partir de \texttt{main}:

\begin{itemize}
  \item \texttt{2863fc4 genes-treinamento: documenta branch de genes treinados do controlador}
  \item \texttt{0f58ccb main: remove genes treinados do versionamento e ignora diretorio}
  \item \texttt{501ce86 artigo: corrige frase sobre rede treinada nao versionada apos remocao do
  versionamento}
\end{itemize}

\texttt{git show 501ce86 -- Artigo/Texto/Artigo.tex} confirma que a frase nova é exatamente a
recomendada pela 3ª revisão, palavra por palavra. A redação está coerente com o estado do repositório
\textbf{em \texttt{main}}: lá, \texttt{git ls-tree -r main} não lista nenhum arquivo em
\texttt{Resources/ControladorTreinado/} além dos \texttt{.cs} de infraestrutura, e
\texttt{git show main:.gitignore} tem a linha \texttt{src/RadarTorres.App/Resources/\allowbreak
ControladorTreinado/}. Um clone de \texttt{main} compilado hoje realmente usa a heurística pura.
\textbf{Até aqui, a correção é exemplar.}

\textbf{O que esta auditoria encontra na branch de trabalho.} O repositório local está em
\texttt{banco-dados} (\texttt{git status}: ``On branch banco-dados''), e esta é a branch cujo commit
de ponta se chama \texttt{a8a7b16 ``revisão do Gabriel no dia da entrega''}, às 14:26 do mesmo dia ---
\textbf{antes} dos dois commits de correção em \texttt{main} (15:03 e 15:40). Nela:

\begin{enumerate}
  \item \texttt{git ls-tree -r banco-dados --name-only} lista \texttt{genoma-controlador.json} e
  \texttt{rede-controlador.net} normalmente, junto dos \texttt{.cs}.
  \item \texttt{git show banco-dados:.gitignore} \textbf{não} tem a linha que exclui o diretório --- a
  busca por ``Controlador'' ou ``genes'' no arquivo não retorna nada.
  \item O \texttt{Artigo.tex} efetivamente no disco (a mesma branch) ainda tem, nas linhas 476--478:
  ``O repositório não versiona uma rede treinada: um executável compilado sem treinamento prévio na
  máquina usa a heurística pura, comportamento de fábrica previsto no projeto.'' --- idêntica à frase
  que a 1ª, a 2ª e a 3ª revisão apontaram como falsa.
  \item O \texttt{.csproj} nesta branch é idêntico ao de \texttt{main} no trecho relevante
  (\texttt{git diff main banco-dados -- RadarTorres.App.csproj} não mostra diferença): ainda embute
  \texttt{rede-controlador.net} como \texttt{EmbeddedResource} sob
  \texttt{Condition="Exists(...)"}, e o arquivo \textbf{existe no disco} (\texttt{Test-Path} confirma
  os dois arquivos). A condição está satisfeita.
\end{enumerate}

\textbf{Logo:} quem compila a partir de \texttt{banco-dados} hoje obtém o mesmo executável com rede
embutida que a 1ª revisão descreveu, e o genoma embutido continua sendo o mesmo que, na 3ª revisão,
não correspondia a nenhum dos dez campeões relatados na Tabela 1 --- essa parte da auditoria não foi
refeita nesta rodada porque o arquivo não mudou, mas nada no repositório sustenta que ele tenha
passado a corresponder a algum resultado reportado.

\textbf{Por que isto é crítico e não apenas um detalhe de processo.} O critério desta banca é
``a afirmação central é falsa ou não sustentada no artefato que a banca vai examinar''. A banca
examina o que está aberto no repositório na hora da defesa. Se for \texttt{banco-dados} --- e todo
sinal disponível aponta para isso: é a branch local, é a que tem o commit nomeado ``revisão... no dia
da entrega'', é a que contém o conteúdo que esta própria banca já elogiou e que \texttt{main} não tem
--- a frase é falsa e o sistema entregue executa uma rede não documentada, exatamente como nas duas
primeiras revisões. O fato de existir uma correção impecável em outra branch não altera o que está no
exemplar examinado, e não há, nesta auditoria, nenhuma evidência de que alguém tenha decidido qual
branch é a oficial.

\textbf{O que não é achado aqui.} Não é acusação de má-fé nem de trabalho perdido: a correção em
\texttt{main} está correta e bem escrita, e o \texttt{README} da branch \texttt{genes-treinamento}
(não reconferido nesta rodada, sem indício de mudança desde a 3ª revisão) continua documentando uma
prática legítima. O achado é específico: \textbf{falta reconciliar as duas branches antes da defesa},
e isso é trabalho de minutos, não de reexperimento.

\textbf{Correção.} Decidir qual branch será apresentada e, se for \texttt{banco-dados} (ou qualquer
outra que não seja \texttt{main}), fazer merge ou cherry-pick dos commits \texttt{0f58ccb} e
\texttt{501ce86} nela, conferindo depois que \texttt{git ls-tree} não lista mais os dois arquivos e
que a Seção 2.3.5 tem a frase nova. Se a intenção for o contrário --- consolidar tudo em
\texttt{main} ---, então os commits exclusivos de \texttt{banco-dados} (explicação da tanh, dívida
técnica dos ciclos de dependência, item (vii), \texttt{cc11e44}, \texttt{a8a7b16}) precisam ser
trazidos para lá, porque hoje \texttt{main} não os tem. De um jeito ou de outro, o estado final
precisa ser uma única branch com as duas coisas.

\subsection*{D.2 \alto\ ``Parâmetros padrão do sistema (população de 60)'' já não são os padrões do
sistema --- persiste}

Reconferido nesta rodada contra o \texttt{Artigo.tex} em disco (linha 235): a frase é idêntica à da
3ª revisão. Não houve alteração de código nesta janela que mude o veredicto. Mantido o diagnóstico e a
correção da 3ª revisão: citar os valores do JSON do experimento (população 60, 10 espécies) em vez de
delegar ao ``padrão do sistema'', que no código é hoje 100.

\subsection*{D.3 \retirado\ Contagem de tabelas do banco de dados}

Mantido retirado, sem reavaliação nesta rodada: o argumento de escopo dos autores, aceito na 3ª
revisão (a Seção 2.5 documenta o núcleo do modelo de dados, não o inventário), continua válido e
nenhum fato novo o contesta.

\subsection*{D.4 \baixo\ Contagem interna inconsistente na mesma frase --- persiste}

Linha 522 do \texttt{Artigo.tex} atual: ``sete cadastrais e de auditoria, e quatro que persistem o
histórico do treinamento'' continua não batendo com as 8+3 linhas da tabela abaixo. Idêntico à 3ª
revisão.

\subsection*{D.5 \alto\ A Tabela 2 não é reproduzível a partir do banco --- persiste}

Linhas 643--644: \textit{fitness} médio 9,11 $\to$ 134,17; melhor por bloco 119,62 $\to$ 137,35.
Nenhum artefato novo foi apresentado nesta rodada para sustentar esses quatro números, e nenhuma
mudança no banco foi relatada. O diagnóstico e a correção da 3ª revisão (exportar e versionar um dump
congelado, citá-lo na legenda) seguem valendo sem alteração.

\subsection*{D.6 \moderado\ ``Soma 10 pontos por alvo interceptado'' --- persiste}

Linha 470: a frase ainda descreve a recompensa como fixa. O código
(\texttt{AmbienteTreinamento.cs:800}, não reaberto nesta rodada por falta de indício de mudança)
soma \texttt{RecompensaAbate * (0.7 + 0.3 * bonusDistancia)}, isto é, 7 a 10 pontos. Correção
inalterada da 3ª revisão.

\subsection*{D.7 \moderado\ ``Cada quadrante admite até cinco alvos simultâneos'' --- persiste}

Linha 469: frase idêntica. Sem teto por quadrante no código, pelo levantamento da 3ª revisão, não
refeito nesta rodada por ausência de qualquer commit que toque \texttt{AmbienteTreinamento.cs}.

\subsection*{D.8 \baixo\ ``4631 oportunidades de acionamento'' --- persiste}

Linha 607: idêntica. Correção inalterada: acrescentar ``em cada execução''.

\subsection*{D.9 \baixo\ Execuções a partir de cópia fora do repositório --- não reconferido}

Não há indício de que os artefatos de resultado tenham sido regenerados nesta janela; mantido como
estava, sem nova verificação do caminho \texttt{C:\textbackslash TEMP} nos JSONs.

\subsection*{D.10 \baixo\ Hardware autodeclarado --- persiste}

Sem mudança: as especificações continuam autodeclaradas e a seção de ambiente computacional não foi
tocada nesta janela.

\section*{E. Auditoria de dados e tabelas}

\begin{longtable}{|p{3.6cm}|p{6.6cm}|p{4.4cm}|}
\hline
\textbf{Elemento} & \textbf{O que foi conferido} & \textbf{Veredicto} \\
\hline
\endhead
Commits \texttt{2863fc4}, \texttt{0f58ccb} (remoção dos genes do versionamento de \texttt{main}) &
\texttt{git log --all}, \texttt{git show}, \texttt{git ls-tree -r main}, \texttt{git show
main:.gitignore} & \textbf{Confere em \texttt{main}} --- diretório não tem arquivo versionado,
\texttt{.gitignore} cobre o caminho \\
\hline
Commit \texttt{501ce86} (reescrita da frase da Seção 2.3.5) & \texttt{git show 501ce86 --
Artigo/Texto/Artigo.tex}: diff compara frase antiga e nova & \textbf{Confere em \texttt{main}} ---
redação idêntica à recomendada na 3ª revisão \\
\hline
Estado de \texttt{banco-dados} (branch de trabalho, commit de ponta \texttt{a8a7b16}) &
\texttt{git ls-tree -r banco-dados}, \texttt{git show banco-dados:.gitignore}, grep no
\texttt{Artigo.tex} em disco, \texttt{git diff main banco-dados -- *.csproj}, \texttt{Test-Path} dos
dois arquivos & \textbf{Correção ausente} --- genoma versionado e embutido, frase antiga, arquivos
presentes no disco \\
\hline
Relação entre \texttt{main} e \texttt{banco-dados} & \texttt{git merge-base}, \texttt{git
merge-base --is-ancestor} (nos dois sentidos), \texttt{git log main..banco-dados},
\texttt{git log banco-dados..main} & \textbf{Divergem sem reconciliação} --- 2 commits exclusivos de
cada lado desde \texttt{b0af938} \\
\hline
Tabela 1 --- comparação pareada & Não reaberta nesta rodada; nenhum artefato ou linha de código que a
produz foi tocado desde a 3ª revisão & \textbf{Sem indício de mudança} --- mantido o veredicto da 3ª
revisão (confere exatamente) \\
\hline
Tabela 2 --- certeza estática & Linhas 643--644 do \texttt{Artigo.tex} reconferidas; banco não
reaberto nesta rodada por ausência de indício de atualização & \textbf{Não reproduz} --- mantido (ver
D.5) \\
\hline
Genoma de produção (branch \texttt{banco-dados}) & Arquivo presente no disco e rastreado pelo git;
conteúdo não recomparado byte a byte nesta rodada por não ter havido alteração desde a 3ª revisão &
\textbf{Mesma divergência da 3ª revisão, agora outra vez relevante} --- não corresponde a nenhum dos
dez campeões relatados (ver D.1) \\
\hline
Conteúdo exclusivo de \texttt{banco-dados} (tanh, dívida técnica, item vii) & \texttt{git diff
banco-dados main -- Artigo/Texto/Artigo.tex} & \textbf{Ausente em \texttt{main}} --- confirma que
nenhuma das duas branches contém a versão completa do artigo \\
\hline
\end{longtable}

\section*{F. Referências}

Nenhuma referência foi adicionada, removida ou reconferida nesta rodada. O julgamento da 3ª revisão
permanece válido para a branch auditada: 21 referências, disciplina de amarrar cada uma a uma decisão
de projeto preservada, lacuna do modelo de ruído ultrassônico ainda sem fonte e declarada como tal.
\textbf{Não verificadas nesta rodada}: todas as 21 --- duas foram verificadas na 1ª revisão; as
demais seguem não verificadas, o que não é o mesmo que corretas.

\section*{G. Pontos fortes}

\begin{enumerate}
  \item \textbf{A correção de D.1, onde foi aplicada, está correta.} Commits \texttt{2863fc4} e
  \texttt{0f58ccb} seguem exatamente o caminho recomendado --- \texttt{git rm --cached} preservando o
  arquivo em disco e numa branch dedicada, com \texttt{.gitignore} e, pelo relato, \texttt{README} de
  propósito --- e o commit \texttt{501ce86} usa a redação exata que a revisão anterior sugeriu, sem
  tentar diluir a afirmação.
  \item \textbf{A tabela central de resultados continua intocada e, até onde esta rodada verificou,
  reproduzível} --- nenhum achado novo a afeta.
  \item \textbf{O resultado negativo permanece declarado no Resumo, Abstract e Conclusão}, sem
  indício de reversão nesta janela.
  \item \textbf{Os oito achados persistentes são todos de correção barata e nenhum exige
  reexperimento} --- o perfil de pendências continua favorável a uma defesa bem-sucedida, uma vez
  resolvida a questão de branch.
\end{enumerate}

\section*{H. Nota final}

\begin{center}
\Large\textbf{7,7 / 10} \quad\normalsize (pareceres anteriores: 6,33; 6,20; 7,9)
\end{center}

A nota caiu 0,2 em relação à 3ª revisão porque o crédito dado então a Reprodutibilidade dependia de a
correção de D.1 estar refletida na branch que a banca vai examinar, e esta rodada encontra que não
está. Resolvida a reconciliação de branches --- sem precisar de nenhum experimento novo --- e
mantidas as demais sete pendências pelo preço que a 3ª revisão já estimou, este artigo volta à faixa
de 8,3--8,5 que a revisão anterior projetou. A distância entre a nota atual e essa faixa não é
método, não é relevância, e não é nem mais a redação de uma frase: é decidir, de forma auditável, qual
branch é o artigo.

\section*{I. Pendências prioritárias}

\subsection*{Bloqueiam a defesa}

\begin{enumerate}
  \item Decidir qual branch será apresentada à banca e levar a ela tanto a correção de D.1
  (\texttt{0f58ccb} + \texttt{501ce86}, hoje só em \texttt{main}) quanto o conteúdo que só existe em
  \texttt{banco-dados} (explicação da tanh, dívida técnica dos ciclos de dependência, item vii).
  Verificar ao final, na branch escolhida: \texttt{git ls-tree -r <branch> --name-only | findstr
  Controlador} não deve listar \texttt{.net} nem \texttt{.json}, e a Seção 2.3.5 deve conter ``a
  versão entregue não inclui rede treinada''. \emph{(D.1 --- é um merge ou cherry-pick de dois
  commits pequenos; nenhum reexperimento.)}
  \item Exportar e versionar um dump congelado com as linhas que produziram a Tabela 2 e citá-lo na
  legenda; se o artefato não for recuperável, declarar isso. \emph{(D.5 --- sem alteração desde a 3ª
  revisão; continua a única afirmação numérica do artigo sem fonte conferível.)}
\end{enumerate}

\subsection*{Alto impacto, esforço baixo}

\begin{enumerate}\setcounter{enumi}{2}
  \item Substituir ``os parâmetros padrão do sistema (população de 60...)'' pelos valores fixos do
  JSON do experimento. \emph{(D.2)}
  \item Trocar ``soma 10 pontos por alvo interceptado'' por ``até 10 pontos (7 a 10, conforme a
  distância)''. \emph{(D.6)}
  \item Remover ``cada quadrante admite até cinco alvos simultâneos'' ou implementar o teto e
  reexecutar. \emph{(D.7)}
\end{enumerate}

\subsection*{Correções de texto}

\begin{enumerate}\setcounter{enumi}{5}
  \item Corrigir a partição ``sete cadastrais e quatro de treinamento'' para bater com as linhas da
  própria tabela (8+3). \emph{(D.4)}
  \item Acrescentar ``em cada execução'' às 4631 oportunidades de acionamento. \emph{(D.8)}
  \item Dizer, na seção de ambiente computacional, que as execuções reportadas partiram de uma cópia
  do projeto fora do repositório. \emph{(D.9)}
  \item Dizer que as especificações de hardware foram coletadas na própria máquina de execução.
  \emph{(D.10)}
\end{enumerate}

\section*{J. Escopo e limitações desta auditoria}

Nenhum arquivo do projeto foi alterado pela banca na produção deste parecer. Todos os comandos usados
para esta rodada foram de leitura: \texttt{git log}, \texttt{git show}, \texttt{git diff}, \texttt{git
ls-tree}, \texttt{git branch}, \texttt{git merge-base}, \texttt{git status}, grep no \texttt{.tex} em
disco e \texttt{Test-Path} nos dois arquivos de rede. Nenhum deles escreve no repositório. Nenhuma
migração ou experimento foi executado.

Esta rodada \textbf{não} refez a reprodução completa da Tabela 1 nem a varredura de referências feita
na 3ª revisão; reconferiu apenas os oito achados persistentes contra o texto em disco (comparação
textual direta, não recálculo numérico) e investigou a fundo apenas D.1, por ser o objeto declarado
desta rodada. Os achados D.2, D.4, D.6--D.10 foram validados por grep de linha, não por reexecução de
script ou reabertura de banco --- coerente com o fato de nenhum commit nesta janela ter tocado o
código ou os dados que os sustentam.

\textbf{Não verificado nesta rodada:}

\begin{itemize}
  \item \textbf{O README da branch \texttt{genes-treinamento}} que a 3ª revisão descreveu; não foi
  reaberto, por ausência de qualquer commit novo nessa branch.
  \item \textbf{Recomparação byte a byte do genoma em disco} contra os dez campeões relatados; mantém-se
  o veredicto da 3ª revisão (não corresponde a nenhum) porque o arquivo não mudou, mas não foi refeita.
  \item \textbf{As fontes bibliográficas.} Seguem não verificadas, à exceção das duas confirmadas na
  1ª revisão.
  \item \textbf{A origem dos números da Tabela 2}, o módulo de visão computacional (declarado não
  avaliado pelo próprio artigo) e o firmware do Arduino.
  \item \textbf{Qual branch será efetivamente apresentada na defesa.} Esta é uma decisão dos autores,
  não algo que o repositório permita inferir com certeza --- a inferência de que \texttt{banco-dados}
  é a branch de entrega se apoia no nome do último commit e no fato de ser a branch localmente
  ativa, não em nenhuma declaração explícita no projeto.
\end{itemize}

\textbf{Artefatos consultados nesta rodada:} \texttt{Artigo/Texto/Artigo.tex} (em disco, branch
\texttt{banco-dados}); histórico git completo (\texttt{git log --all}, \texttt{--oneline}, com datas);
\texttt{git show} dos commits \texttt{2863fc4}, \texttt{0f58ccb}, \texttt{501ce86}, \texttt{a8a7b16};
\texttt{git ls-tree -r} de \texttt{main} e \texttt{banco-dados}; \texttt{.gitignore} de ambas as
branches; \texttt{git diff main banco-dados} sobre \texttt{Artigo.tex} e
\texttt{RadarTorres.App.csproj}; \texttt{src/RadarTorres.App/RadarTorres.App.csproj};
\texttt{src/RadarTorres.App/Resources/ControladorTreinado/\{rede-controlador.net,
genoma-controlador.json\}} (existência em disco); \texttt{Artigo/Revisoes/Parecer\_Banca\_TCC\_rev3.tex}
(parecer anterior, íntegro).

\end{document}
